\section{Limitations}

\paragraph{Source data and task design.}
The 142 source pairs are drawn from English OASST2, so the experiment measures
controlled translation and romanization rather than naturalistic response to
real romanized user input. Users who type in romanized script produce
organically code-mixed, spelling-variable queries that differ systematically
from the formal English prompts used here. Results may underestimate real-world
quality degradation for models that struggle with informal romanized input.

\paragraph{Single prompting strategy.}
All four models receive the same structured JSON prompt. Different prompting
strategies---chain-of-thought, few-shot examples, native-script intermediate
steps---may substantially change output quality, particularly for Ollama's
Kannada and Telugu failures. The results characterize one prompting regime,
not the ceiling of each model's romanization capability.

\paragraph{Automatic metrics.}
BGE-M3 cosine similarity is our primary metric, but it assigns passing scores
to verbatim English copies (BGE = 1.000) and native-script outputs
(BGE $\approx$ 0.80--0.84), demonstrating its own limitations for this task.
Raw BERTScore F1 using \texttt{bert-base-multilingual-cased} is difficult to
interpret for cross-lingual English-to-romanized comparison, as its token
alignment was not designed for this setting. Neither metric captures
naturalness, register, or spelling consistency.

\paragraph{Human evaluation coverage.}
Human ratings are collected on a purposively sampled subset of approximately 24 responses
per language (73 total) rather than the full 142-row dataset. Evaluators are sampled from
a single team; ratings from speakers across different regions, age groups, and
levels of digital literacy may differ. Romanization has no single universally
correct spelling, and rater agreement on naturalness judgments was not
formally measured. The human evaluation results should be interpreted as
directional evidence rather than a definitive quality ranking.

\paragraph{Languages and models.}
We study Hindi, Telugu, and Kannada. Other Indic romanized varieties---Bengali,
Tamil, Marathi---may exhibit different patterns. The four models evaluated here
(Claude, GPT-4o, Gemini, Qwen3.5-35B) represent a snapshot of current capabilities; results for other model families or future versions may differ.
